% §5.4 Lessons learned -- replaces the section of the same name in the draft of 10 Nov 2026.
% Airports: the 108 tasks resampled with the new difficulty measure (data/new-diff-airports_ressampled.csv);
% football: the full context-on task set (239 tasks). Figures: notebooks/figures/paper_reports/airports_resampled/*.pdf
% (copy the folder to figures/airports_resampled/). Needs \usepackage{booktabs} and the macro
% \newcommand{\fign}{F1\textsubscript{ign}}. Generated by notebooks/make_lessons_learned_tex.py.
% Significance: 95% paired bootstrap over tasks (2,000 resamples).
% Labels expected in the draft: sec:difficulty (Section 3.2, Task Difficulty), sec:effort (Section 4.3.2, Effort) and
% sec:eis (Section 4.3.3, energy).
% As requested in the draft (Patrick, Vero), only the new results are reported, with one closing paragraph on DOLAP.

\subsection{Lessons learned}
\label{sec:lessons}

This section reports the accuracy of both strategies, the effect of context and of task difficulty, and the cost of
each strategy in energy and effort. On airports, all results use the 108 tasks sampled with the global difficulty measure
of Section~\ref{sec:difficulty}; on football, the 239 tasks of the context-on set. Each configuration is run three times and
scored on its voted answer. We report F1 and precision (P) as the main measures and \fign{} for comparison with the DOLAP
study. \emph{A mean} denotes the average over the nine prompts of Strategy~A and \emph{A best} the best prompt of each model
by F1, chosen on the same tasks, which makes it an optimistic reference.

\subsubsection{Strategy A: the LLM alone}

\paragraph{Main results.}
On airports, Strategy~A reaches F1 0.20--0.25 averaged over the prompts and 0.22--0.31 with the best prompt of each model
(Table~\ref{tab:strategies}); precision is 0.27--0.42 and 0.31--0.50 respectively. Between 12\% and 57\% of the hidden
properties are left unanswered with the best prompt, and 21--59\% are answered wrongly. On football, the best prompt reaches
F1 0.19--0.43 and precision 0.25--0.49, with deepseek-v4.1-flash and mistral-large-3 ahead of the other three models.

\paragraph{Model or prompt?}
No prompt is best for every model (Table~\ref{tab:a-prompts}): the chain-of-thought variants are best for four models on
airports and the compact prompt for gpt-oss:120b. On airports, the choice of prompt matters as much as the choice of model:
for a fixed model the nine prompts span 0.11 F1 on average, and for a fixed prompt the five models span 0.09. Averaged over
the models, however, every prompt scores F1 0.22--0.25, so no prompt is better in general. On football the model dominates:
models span 0.28 F1 for a fixed prompt, prompts 0.16 for a fixed model, and the compact prompt is best for three models.

\begin{table}[t]
\centering
\caption{Airports (108 tasks): Strategy A: F1 and precision (P) of each model and prompt, voted over 3 runs. Bold: each model's best prompt by F1 (A best).}
\label{tab:a-prompts}
\small
\setlength{\tabcolsep}{3.5pt}
\begin{tabular}{lrrrrrrrrrrr}
\toprule
 & \multicolumn{2}{c}{gpt-oss:20b} & \multicolumn{2}{c}{gemma4:31b} & \multicolumn{2}{c}{gpt-oss:120b} & \multicolumn{2}{c}{mistral-large-3} & \multicolumn{2}{c}{deepseek-v4.1-flash} & \\
\cmidrule(lr){2-3}\cmidrule(lr){4-5}\cmidrule(lr){6-7}\cmidrule(lr){8-9}\cmidrule(lr){10-11}
Prompt & F1 & P & F1 & P & F1 & P & F1 & P & F1 & P & Mean F1 \\
\midrule
compact & 0.21 & 0.22 & 0.24 & 0.30 & \textbf{0.30} & 0.33 & 0.16 & 0.16 & 0.19 & 0.20 & 0.22 \\
zero-shot & 0.15 & 0.21 & 0.26 & 0.44 & 0.25 & 0.42 & 0.24 & 0.28 & 0.21 & 0.27 & 0.22 \\
zero-shot + CoT & 0.21 & 0.27 & \textbf{0.30} & 0.50 & 0.26 & 0.50 & 0.21 & 0.25 & 0.21 & 0.27 & 0.24 \\
1-shot & 0.20 & 0.27 & 0.23 & 0.40 & 0.21 & 0.38 & 0.23 & 0.24 & 0.25 & 0.31 & 0.22 \\
1-shot + CoT & 0.18 & 0.25 & 0.23 & 0.40 & 0.25 & 0.47 & 0.26 & 0.29 & 0.27 & 0.36 & 0.24 \\
3-shot & 0.20 & 0.29 & 0.18 & 0.31 & 0.22 & 0.38 & 0.26 & 0.28 & 0.27 & 0.32 & 0.23 \\
3-shot + CoT & \textbf{0.22} & 0.31 & 0.20 & 0.33 & 0.27 & 0.48 & \textbf{0.29} & 0.33 & 0.23 & 0.27 & 0.24 \\
5-shot & 0.21 & 0.29 & 0.23 & 0.40 & 0.23 & 0.40 & 0.26 & 0.27 & 0.20 & 0.23 & 0.23 \\
5-shot + CoT & 0.22 & 0.30 & 0.23 & 0.39 & 0.20 & 0.39 & 0.26 & 0.29 & \textbf{0.31} & 0.36 & 0.25 \\
\bottomrule
\end{tabular}
\end{table}


\begin{table}[t]
\centering
\caption{Airports (108 tasks): mean and best results of each strategy per model. A mean averages the 9 prompts; A best is the model's best prompt by F1 (Table~\ref{tab:a-prompts}). Energy per run as defined in Section~\ref{sec:eis}; energy per correct answer counts the three runs of every task.}
\label{tab:strategies}
\small
\setlength{\tabcolsep}{4pt}
\begin{tabular}{llrrrrrr}
\toprule
Model & Strategy & F1 & P & \fign & Abst. & Wh/run & Wh/correct \\
\midrule
gpt-oss:20b & A mean & 0.20 & 0.27 & 0.42 & 38\% & 0.009 & 0.16 \\
 & A best & 0.22 & 0.31 & 0.47 & 43\% & 0.009 & 0.16 \\
 & B search & 0.32 & 0.41 & 0.58 & 34\% & 0.707 & 7.90 \\
 & B search + memory & 0.30 & 0.34 & 0.51 & 21\% & 0.708 & 7.91 \\
\addlinespace
gemma4:31b & A mean & 0.23 & 0.39 & 0.55 & 55\% & 0.001 & 0.02 \\
 & A best & 0.30 & 0.50 & 0.67 & 57\% & 0.001 & 0.01 \\
 & B search & 0.21 & 0.37 & 0.54 & 60\% & 0.358 & 7.26 \\
 & B search + memory & 0.22 & 0.32 & 0.48 & 47\% & 0.358 & 6.45 \\
\addlinespace
gpt-oss:120b & A mean & 0.24 & 0.42 & 0.59 & 56\% & 0.023 & 0.39 \\
 & A best & 0.31 & 0.33 & 0.49 & 12\% & 0.035 & 0.36 \\
 & B search & 0.42 & 0.56 & 0.72 & 41\% & 0.975 & 8.78 \\
 & B search + memory & 0.41 & 0.51 & 0.68 & 33\% & 0.976 & 8.55 \\
\addlinespace
mistral-large-3 & A mean & 0.24 & 0.27 & 0.42 & 17\% & 0.092 & 1.26 \\
 & A best & 0.29 & 0.33 & 0.50 & 22\% & 0.129 & 1.49 \\
 & B search & 0.35 & 0.44 & 0.61 & 32\% & 1.692 & 17.13 \\
 & B search + memory & 0.28 & 0.33 & 0.50 & 24\% & 1.790 & 21.48 \\
\addlinespace
deepseek-v4.1-flash & A mean & 0.24 & 0.29 & 0.44 & 28\% & 0.041 & 0.61 \\
 & A best & 0.31 & 0.36 & 0.53 & 26\% & 0.046 & 0.52 \\
 & B search & 0.18 & 0.30 & 0.47 & 57\% & 1.281 & 29.64 \\
 & B search + memory & 0.25 & 0.36 & 0.53 & 46\% & 1.285 & 19.82 \\
\bottomrule
\end{tabular}
\end{table}

\subsubsection{Context}

Figure~\ref{fig:context} compares each model with and without context on the 77 airport tasks that exist in both
conditions. Context lowers the accuracy: pooled over the models and prompts, F1 is 0.35 without context and 0.27 with it,
precision 0.40 and 0.34, and abstention rises from 23\% to 35\%. Context scores higher in only 1 of the 45 model--prompt
pairs. This reverses the finding of the DOLAP study, where context helped on the hardest tasks. The two designs differ: in
the DOLAP study, context entries were added to the node's own properties and 82--88\% of them described neighbouring
airports; here $G$ has a fixed size, so context entries replace some of the node's own properties, and they are mostly
properties of the airport's country (44\% of the entries, against 33\% for neighbouring airports and 22\% for the routes
and passenger flows that connect them).
In 13 of the 108 tasks, $G$ contains no property of the airport at all. The models need to identify the airport, which its
own properties allow and country statistics rarely do.

\begin{figure*}[t]
\centering
\includegraphics[width=\textwidth]{figures/airports_resampled/dolap_fig3_accuracy_vs_difficulty.pdf}
\caption{Airports (77 tasks present with and without context): F1 (top) and precision (bottom) of each model without context ($x$) and with context ($y$), by thirds of the global difficulty, all prompts of Strategy~A pooled. Points above the diagonal benefit from context.}
\label{fig:context}
\end{figure*}

\subsubsection{Accuracy and task difficulty}

The global difficulty measure orders the tasks for every strategy (Figure~\ref{fig:difficulty}). From the easiest to the
hardest third of the tasks, F1 falls from 0.56 to 0.21 and 0.06 with A best, from 0.44 to 0.27 and 0.17 with search, and
from 0.46 to 0.25 and 0.15 with search and memory; precision follows the same order. Per task, the share of correct answers
over the 15 model--strategy combinations decreases with difficulty (Spearman $\rho = -0.55$), more steeply for A best
($-0.64$) than for search ($-0.34$). Of the 108 tasks, 39 are never answered correctly, by any model or strategy, and only 2
by all 15 combinations (Figure~\ref{fig:where}). The most frequent hidden property, \texttt{commercial\_flights} (28 tasks),
is found in 4--7\% of the answers. By comparison, the strata used to build the original task set order the tasks less
consistently: with A best, tasks whose hidden property has a medium popularity score higher (F1 0.42) than those with a
popular one (0.28).

\begin{figure*}[t]
\centering
\includegraphics[width=\textwidth]{figures/airports_resampled/accuracy_vs_difficulty_bins.pdf}
\caption{Airports (108 tasks): F1 (top) and precision (bottom) by thirds of the global difficulty, per model (lines) and pooled (dashed), for A best, B search and B search + memory.}
\label{fig:difficulty}
\end{figure*}

\begin{figure*}[t]
\centering
\includegraphics[width=\textwidth]{figures/airports_resampled/where_models_are_right.pdf}
\caption{Airports (108 tasks): outcome of every task (columns, from the easiest to the hardest) for each model and strategy (rows). Top: difficulty; bottom: number of correct answers out of 15.}
\label{fig:where}
\end{figure*}

\subsubsection{Strategy B: search}

With search, gpt-oss:120b reaches the highest F1 of the study on airports (0.42, precision 0.56), followed by
mistral-large-3 (0.35) and gpt-oss:20b (0.32); gemma4:31b and deepseek-v4.1-flash stay at 0.21 and 0.18, abstaining on
57--60\% of the tasks (Figure~\ref{fig:accuracy}). Against A best, search raises F1 significantly only for gpt-oss:120b
($+0.11$, with $+0.24$ precision) and lowers it significantly for deepseek-v4.1-flash ($-0.13$); for the other three models
the difference is not significant. Pooled over the models, search is the most precise strategy (precision 0.43 against 0.35
for A best) at a similar F1 (0.30 against 0.29): it replaces wrong answers by abstentions (wrong answers fall from 44\% to
32\%, abstentions rise from 32\% to 45\%). Search finds the facts printed on an airport's page better than the model's
memory (\texttt{runwaysWidth}: 48\% against 15\%; \texttt{runwaysLength}: 33\% against 10\%) and loses on the number of
runways (18\% against 62\%).

On football, search lowers F1 for every model, significantly for four of them ($-0.11$ to $-0.36$; gpt-oss:120b $-0.03$, not
significant): the agent abstains on 54--87\% of the properties and reaches F1 0.07--0.17.

\begin{figure*}[t]
\centering
\includegraphics[width=\textwidth]{figures/airports_resampled/accuracy_by_approach.pdf}
\caption{Airports (108 tasks): F1, precision and \fign{} per model for Strategy~A (mean of the prompts, best prompt) and Strategy~B (search, search + memory).}
\label{fig:accuracy}
\end{figure*}

\subsubsection{Strategy B: search + memory}

Allowing the agent to answer from its own knowledge lowers abstention by 7--13 points for every model on airports, but the
additional answers are more often wrong: precision falls for four of the five models. The change in F1 is significant only
for mistral-large-3 ($-0.07$) and deepseek-v4.1-flash ($+0.07$). On football, memory raises F1 for every model (0.12--0.23
against 0.07--0.17 with search alone), but remains below A best for four of the five models; only gpt-oss:120b is slightly
above (0.20 against 0.19).

\subsubsection{The LLM alone against the agent}

The two strategies are complementary rather than ranked. On airports, A best is more accurate on the easiest third of the
tasks (F1 0.56 against 0.44 for search), and search on the medium and hardest thirds (0.27 against 0.21, and 0.17 against
0.06). Of the 69 tasks answered correctly at least once, 19 are answered correctly only with search and 9 only with A best.
On football, the LLM alone is the better strategy for four of the five models; for gpt-oss:120b the two are level.

\subsubsection{Energy consumption}
\label{sec:lessons-energy}

We report the energy consumption of a run, computed with the formula of Section~\ref{sec:eis}: the EcoLogits estimate of
the run's LLM calls plus 0.3\,Wh per web search. We report energy in Wh rather than a composite score: water consumption
ranks the 55 configurations exactly as energy does and global warming potential almost so (Spearman $\rho = 0.87$; it
differs only through the electricity mix EcoLogits assumes for each provider), and with a single indicator a min--max
normalisation would only rescale it relative to the configurations compared.

Strategy~A runs use 0.0007--0.13\,Wh (Figure~\ref{fig:energy}). The model sets the energy budget: gemma4:31b uses about
0.001\,Wh per run and mistral-large-3 about 0.09\,Wh, a 95-fold span, and gpt-oss:20b uses nine times the energy of the
larger gemma4:31b because it still produces reasoning tokens. Within a model, the prompts differ by a factor of 1.5--4.7,
and energy follows the number of output tokens (Spearman $\rho \geq 0.97$) rather than the length of the prompt ($-0.27$
to $0.73$); the compact prompt, the shortest, is the most expensive for four of the five models because it leads to the
longest answers. Across the Strategy~A configurations, energy and F1 are only weakly related ($\rho = 0.25$).

A run of Strategy~B costs 0.36--1.79\,Wh, i.e.\ 13 (mistral-large-3) to 545 (gemma4:31b) times the energy of the model's
best Strategy~A prompt, and the searches account for 53--98\% of it. Per correct answer, counting the three runs of every
task, Strategy~A costs 0.009--1.7\,Wh and Strategy~B 6.4--29.6\,Wh. The most accurate configuration (gpt-oss:120b with
search, F1 0.42) costs 0.98\,Wh per run and 8.8\,Wh per correct answer, whereas gemma4:31b with its best prompt reaches
F1 0.30 at 0.0007\,Wh per run.

\begin{figure*}[t]
\centering
\includegraphics[width=\textwidth]{figures/airports_resampled/energy.pdf}
\caption{Airports (108 tasks): energy per run (left) and per correct answer (middle) of each model and strategy, and F1 against energy per run for every configuration (right); log scales.}
\label{fig:energy}
\end{figure*}

\subsubsection{Effort}
\label{sec:lessons-effort}

We measure effort with the token amplification $|T_{process}|/|T_{user}|$ of Section~\ref{sec:effort}, i.e.\ the tokens a
strategy processes in a run (input and output, reasoning included, over all LLM calls) per token the user writes for the
task (Figure~\ref{fig:amplification}).

\paragraph{Strategy A.}
Strategy~A processes 1.2--3.8 times the tokens written by the user, depending on the model and the prompt. On average over
the prompts, the amplification is 1.31--1.43 for deepseek-v4.1-flash, mistral-large-3 and gemma4:31b, whose answers are a
few tokens long, and 1.87--2.08 for the two gpt-oss models, whose output is 88--92\% reasoning. Across prompts, the
amplification falls as the user writes more, from 2.53 for the compact prompt (135 tokens written per task) to 1.39 for
5-shot with chain of thought (481 tokens): a higher amplification reflects a shorter input more than additional work.
More effort from the user does not buy accuracy: the mean F1 over the five models stays between 0.22 and 0.25 for every
prompt, and across the 45 configurations, amplification and F1 are weakly and negatively related (Spearman
$\rho = -0.30$), with a sign that changes from one model to another ($-0.54$ to $0.58$).

\paragraph{Strategy B.}
For a comparable written input (579 tokens per task), Strategy~B processes 5.8--7.2 times the tokens written by the user
(5.7--7.2 with memory). Most of these tokens are read rather than generated: the retrieved evidence accounts for 65--74\% of
the processed tokens (2,180--2,760 per run), while the output remains shorter than the user's input (0.17--0.74 times its
length). Over the 55 configurations of both strategies, amplification does not predict accuracy ($\rho = 0.05$): the agent
is both the most amplified and the most accurate configuration for the two gpt-oss models and mistral-large-3, but not for
gemma4:31b and deepseek-v4.1-flash. The agent therefore shifts effort from the user to the system, and its accuracy gains, where they
exist, are paid in tokens read and, above all, in web searches (Section~\ref{sec:lessons-energy}).

\begin{figure*}[t]
\centering
\includegraphics[width=\textwidth]{figures/airports_resampled/token_amplification.pdf}
\caption{Airports (108 tasks): token amplification $|T_{process}|/|T_{user}|$ per model and strategy (left), per Strategy~A prompt averaged over the models, with the tokens written per task in brackets (middle), and F1 against amplification for every configuration (right).}
\label{fig:amplification}
\end{figure*}

\subsubsection{Comparison with the DOLAP study}

The scores of the DOLAP study are not directly comparable with ours: they were computed over answered properties only,
including properties that were not asked, and report the best of 18 configurations. Re-scored with the measures of this
paper and pooling all configurations, the six proprietary models of that study reach F1 0.17--0.43, precision 0.25--0.54 and
\fign{} 0.39--0.70 on its 36 tasks. To compare the models directly, we ran the four open-weight models that remain
available (mistral-large-3 was withdrawn from the serving platform) on the same 36 tasks, with the published prompts and five
runs. Pooling all configurations, deepseek-v4.1-flash (F1 0.31) and gpt-oss:120b (0.28) fall between Gemini 2.5 Pro (0.33)
and Gemini 2.5 Flash (0.28), and gemma4:31b (0.20) and gpt-oss:20b (0.18) at the level of the two Claude models
(0.17--0.22); only GPT-5 (0.43) and GPT-5 Mini (0.38) remain clearly ahead. The open-weight models abstain less (10--27\%
of the hidden properties) and are wrong more often, so their precision is lower (0.19--0.32). With the prompts of this paper,
every model gains 0.05--0.12 F1 on the same tasks, and deepseek-v4.1-flash reaches 0.41. Conversely, the same models and
prompts score 0.03--0.17 F1 lower on our 108 airport tasks than on the DOLAP tasks. The lower scores reported in this paper
therefore come mainly from the stricter scoring and the harder task set; only the GPT-5 models remain clearly ahead of the
open-weight models.
% Source: notebooks/dolap2026_lineup_report.ipynb (configs/strategy_a_dolap2026.yaml, strategy_a_dolap2026_newprompts.yaml).
